\scanpage{94}
% L1, centered title
\begin{center}Fonctions elliptiques de Legendre\end{center}

% L2–3
$\operatorname{sn}u$ est fonction inverse de $\int_0^{\unsure{v|u}}\frac{dx}{\sqrt{(1-x^2)(1-k^2x^2)}}$, impaire. On pose

% L4–6
\[
\left.\begin{aligned}
\operatorname{cn}^2u+\operatorname{sn}^2u&=1\\
\operatorname{dn}^2u+k^2\operatorname{sn}^2u&=1\\
(\operatorname{dn}^2u-k^2\operatorname{cn}^2u&=1-k^2=k'^{\,2})
\end{aligned}\right\}
\quad\begin{aligned}\operatorname{sn}0&=0\\\operatorname{cn}0&=\operatorname{dn}0=1\end{aligned}
\tag*{\textcircled{\scriptsize 1}}
\]

% L7
($\operatorname{cn}$ et $\operatorname{dn}$ sont uniformes, paires)

% L8
On a les équations différentielles

% L9–11
\[
\left\{\begin{aligned}
\operatorname{sn}'u&=\sqrt{(1-\operatorname{sn}^2u)(1-k^2\operatorname{sn}^2u)}=\operatorname{cn}u\operatorname{dn}u\\
\operatorname{cn}'u&=\sqrt{(1-\operatorname{cn}^2u)(k'^{\,2}+k^2\operatorname{cn}^2u)}=-\operatorname{sn}u\operatorname{dn}u\\
\operatorname{dn}'u&=\sqrt{(1-\operatorname{dn}^2u)(\operatorname{dn}^2u-k'^{\,2})}=-k^2\operatorname{sn}u\operatorname{cn}u
\end{aligned}\right.
\tag*{\textcircled{\scriptsize 2}}
\]

% L12
Les développements de Taylor sont, en $0$

% L13–15
\[
\left\{\begin{aligned}
\operatorname{sn}u&=u-\frac{2\alpha k}{3!}u^3+\frac{4k^2(\alpha^2+3)}{5!}u^5-\frac{8k^3(\alpha^3+33\alpha)}{7!}u^7+\ldots\\
\operatorname{cn}u&=1-\frac{u^2}{2!}+\frac{1+4k^2}{4!}u^4-\frac{1+44k^2+16k^4}{6!}u^6+\text{---}\\
\operatorname{dn}u&=1-\frac{k^2u^2}{2!}+\frac{k^2(4+k^2)}{4!}u^4-k^2\frac{16+44k^2+k^4}{6!}u^6+\ldots
\end{aligned}\right.
\tag*{\textcircled{\scriptsize 3}}
\]

% L16
\[\left(\alpha=\frac12\left(k+\frac1k\right)\right)\]

% L17–18
\[
\text{\unsure{Dégénérescences|Dégénérescence}}\quad
\left\{\begin{aligned}
k=0&\quad\operatorname{sn}=\sin&\operatorname{cn}&=\cos,&\operatorname{dn}&=1\\
k=1&\quad\operatorname{sn}=\operatorname{th}&\operatorname{cn}&=\frac1{\operatorname{ch}}&\operatorname{dn}&=\frac1{\operatorname{th}}
\end{aligned}\right.
\]

% L19–20, double underline
\uuline{Périodes et valeurs remarquables}\quad Posons

% L19–20, right of heading
\[
K=\int_0^1\frac{\unsure{du|dx}}{\sqrt{(1-x^2)(1-k^2x^2)}}\qquad
K'=\int_0^1\frac{dx}{\sqrt{(1-x^2)(1-k'^{\,2}x^2)}}
\]

% L21–22, foot
$\operatorname{sn}u,\ldots$ ont les périodes $4K$ et $4iK'$. Posons $K''=-\unsure{k|K}-iK'$, ($4K''$ est une autre période)

% lower left
% archive sheet number: 93
